
\chapter{3-D Secure}\label{ch:3ds}
В~операциях без физического предъявления карты (\textit{card-not-present}, CNP) продавец
оценивает риск до~авторизации и~старается не~превратить оформление заказа
в~долгую проверку, из-за которой покупатель прервёт оформление. 3-D Secure
разрешает это противоречие: эмитент проверяет держателя карты,
а~правила сети определяют, как при~этом перераспределяется ответственность
за~фрод.

\section{Происхождение 3DS: фрод по CNP}\label{sec:3ds-cnp-fraud}

В~базовой модели CNP продавец располагает лишь карточными реквизитами: номером
карты, сроком действия, CVV2 и, возможно, адресом или почтовым индексом. Этих
данных достаточно, чтобы сформировать авторизационный запрос, но недостаточно,
чтобы доказать, что покупку инициирует именно законный держатель карты.
Реквизиты можно выманить фишингом или получить из~утечки, и~после этого
злоумышленник получает возможность инициировать CNP-платежи как реальный клиент.

По~мере того как чип EMV сделал клонирование физической карты менее выгодным,
атаки сместились в~дистанционный канал. Миграция на~чип в~точке продажи заняла
больше десяти лет, и~всё это время канал CNP оставался без~сопоставимой
защиты~--- мошенники перенесли атаки туда.

На~европейском рынке это подтверждают данные ЕЦБ. В~отчёте по~SEPA фрод по~CNP
составил около 84\,\% стоимости карточного мошенничества и~в~2020,
и~в~2021~году; этот сдвиг ЕЦБ объясняет ростом электронной коммерции и~усилением
защиты канала с~физическим предъявлением карты~\cite{ecb-card-fraud-2021}. Даже
на~фоне снижения абсолютной стоимости фрода по~CNP в~2021~году этот канал
оставался основным источником карточных потерь у~европейских эмитентов.

\highlight{Для транзакций CNP карточным сетям был нужен механизм, который решал
  бы ту же задачу, что чип решает в~терминале: давал независимое подтверждение
того, что транзакцию инициировал законный держатель карты. Им стал 3-D Secure}.
Сначала рынок знал этот протокол по~брендам \enquote{Verified by Visa}
и~\enquote{Mastercard SecureCode}, а~затем по~единой спецификации
EMVCo~\cite{emvco-3ds-whitepaper}.

Распространённость 3DS по~рынкам остаётся неоднородной. В~Европе
обязательность усиленной аутентификации плательщика (Strong Customer
Authentication, SCA) по~второй европейской директиве о~платёжных услугах
(Payment Services Directive~2, PSD2~\cite{psd2}; см.~раздел~\ref{sec:psd2-sca})
привела к~массовому внедрению 3DS~2; в~США обязательного требования нет,
и~продавцы балансируют конверсию и~защиту. В~России
3DS используется повсеместно по~инициативе эмитентов и~по~правилам системы Мир.

\begin{figure}[tbp]
  \centering
  \includefiguresvg[width=\linewidth]{assets/figures/ch11-3ds/cnp-attack-surface}
  \caption{Пять векторов компрометации CNP и~зона эффективности 3DS.}%
  \label{fig:cnp-attack-surface}
\end{figure}

Magecart~--- общее имя семейства группировок (Group~1--12+), ведущих
JavaScript-скимминг с~2015~года: вредоносный код внедряется в~CMS магазина или
приходит через сторонние теги (аналитика, чаты, A/B-тестирование) и~встраивается
в~поля ввода на~странице оплаты. AVS (Address Verification Service) сверяет номер дома
и~ZIP и~работает только в~США, Канаде и~Великобритании; в~остальных регионах
аутентификацию держателя карты выполняет 3DS. 3DS снижает риск фишинга
и~ATO, частично помогает против Magecart и~перебора краденых карт (\textit{card testing}), но не~помогает против
\textit{fullz}/AVS. После
первой авторизации PAN на~стороне PSP заменяется сетевым токеном
(глава~\ref{ch:tokenization}), и~уже токен, а~не~открытый PAN, хранится
в~\textit{card-on-file} (CoF~--- сохранённые продавцом реквизиты карты для последующих
списаний) и~используется в~последующих операциях.

\section{Три домена 3DS}\label{sec:3ds-three-domains}

Название протокола, 3-D Secure, означает \enquote{безопасность в~трёх
доменах}. В~спецификации EMVCo они называются \textit{Issuer Domain},
\textit{Acquirer Domain} и~\textit{Interoperability Domain}; далее используются
краткие имена~--- домен эмитента, домен эквайера и~домен интеграции
соответственно. В~домене эмитента находится ACS (Access Control Server)~---
сервер эмитента или его процессора, который принимает решение об~аутентификации
держателя карты. В~домене эквайера со~стороны продавца различают \textit{3DS
Requestor}~--- роль инициатора аутентификации (продавец или его
PSP~--- Payment Service Provider, платёжный провайдер) и~\textit{3DS
Server}~--- технический компонент, который от~имени Requestor формирует
Authentication Request (AReq) и~общается с~Directory Server. Если продавец
не~поддерживает собственную 3DS-инфраструктуру, PSP выполняет роль Requestor,
предоставляет 3DS Server и~встраивает его в~инфраструктуру продавца. Между
доменами эмитента и~эквайера расположен домен интеграции, в~котором Directory
Server карточной сети определяет по~данным карты нужный ACS и~доставляет
сообщения между сторонами.

Распределение ролей жёсткое: 3DS Server передаёт контекст, DS маршрутизирует,
ACS принимает решение. Продавец клиента не~подтверждает; сеть в~решение
эмитента не~вмешивается. От~v1 к~v2 менялись интерфейсы, объём передаваемых
данных и~способы подтверждения, но это разделение оставалось прежним.

Третий домен, домен интеграции, необходим, потому что эквайер и~эмитент
не~знают друг друга напрямую. Эквайер не~может по~одному PAN найти нужный ACS:
у~сотен эмитентов разные реализации, разные эндпоинты и~разные версии
протокола. Directory Server маршрутизирует запрос: по~BIN карты (Bank
Identification Number, банковский идентификационный диапазон) определяет нужный
ACS и~доставляет ему сообщение. Без~этого промежуточного звена каждый эквайер
был~бы вынужден хранить адресную книгу всех эмитентов и~обновлять её при~каждом
изменении адресов ACS.

\begin{figure}[htbp]
  \centering
  \includefiguresvg[width=\linewidth]{assets/figures/ch11-3ds/3ds-three-domains}
  \caption{Архитектура трёх доменов 3-D Secure.}\label{fig:3ds-three-domains}
\end{figure}

\section{3DS v1: редирект и~его проблемы}\label{sec:3ds-v1}

Первая версия протокола работала через браузерный редирект. При~оплате на~сайте
продавца браузер перенаправлялся на~страницу ACS эмитента, где держатель карты
вводил пароль, связанный с~программой \enquote{Verified by Visa},
\enquote{SecureCode} или её локальным аналогом. На~стороне продавца за~обмен
с~Directory Server отвечал MPI (Merchant Plug-In)~--- встроенный компонент 3DS
v1, формировавший пару VEReq/VERes (Verify Enrollment Request/Response) для
проверки участия карты в~программе и~пару PAReq/PARes (Payer Authentication
Request/Response) для собственно аутентификации.

Опыт держателя карты зависел от~статического пароля, который требовался раз
в~несколько месяцев. Его путали с~PIN и~с~паролем от~интернет-банка,
вводили на~странице с~незнакомым доменом~--- и~часть покупателей прерывала
оформление. Переход на~отдельный экран подтверждения показывал клиенту, что
оплата происходит вне магазина, и~доверие к~покупке
падало~\cite{visa-secure-ui-3ds1-0-2}. К~концу жизненного цикла 3DS v1 многие
эмитенты заменили статический пароль на~одноразовый код по~SMS, но сценарий
полноэкранного редиректа на~чужой домен и~связанные с~ним риски остались
прежними.

Редирект на~незнакомый домен упрощал фишинг: пользователь привыкал
вводить пароль на~странице, не~похожей ни~на~магазин, ни~на~привычный
интернет-банк, и~атакующему оставалось воспроизвести похожий экран ACS
на~подконтрольном домене.

С~переходом электронной коммерции на~мобильные устройства проблема обострилась.
3DS v1 был спроектирован для~десктопного браузера: редирект посреди оформления
заказа ещё был приемлем на~большом экране, но в~мобильном вебе и~тем
более в~нативных приложениях он разрывал сценарий покупки.

Логика принятия решения была грубой. ACS в~3DS v1 выбирал одну
из~двух крайностей: пропустить транзакцию без~явной аутентификации или
потребовать статический пароль. У~него не~было ни~достаточного контекста
покупки, ни~встроенного способа перенести подтверждение в~мобильное приложение
банка. Первая версия решала экономическую задачу переноса ответственности, но
плохо решала задачу точной и~удобной аутентификации.

При всех недостатках эмитенты поддерживали 3DS v1, потому что протокол переносил
ответственность (\textit{liability shift}): если продавец использовал 3DS
и~аутентификация была одобрена, риск мошенничества смещался на~эмитента,
а~продавец получал защиту от~chargeback по~причине фрода. Для продавцов это был
компромисс: потеря конверсии в~обмен на~снижение доли chargeback, а~не однозначная
выгода~\cite{emvco-3ds-whitepaper}. Карточные сети прекратили поддержку 3DS v1
в~октябре 2022~года: с~этого момента аутентификация в~электронной коммерции целиком
перешла на~EMV 3DS~v2~\cite{visa-3ds1-sunset-2022}.

\begin{figure}[htbp]
  \centering
  \includefiguresvg[width=\linewidth]{assets/figures/ch11-3ds/3ds-v1-flow}
  \caption{Поток 3DS v1: VEReq/VERes между MPI и~DS, PAReq/PARes через браузер.}\label{fig:3ds-v1-flow}
\end{figure}

\FloatBarrier
\section{3DS v2: аутентификация на~основе риска}\label{sec:3ds-v2}

Вторая версия 3DS сохранила дополнительную аутентификацию, но отказалась
от~обязательного редиректа и~статического пароля. EMVCo развивает это семейство
спецификаций как набор версий 2.x~\cite{emvco-3ds-whitepaper}. Главное
нововведение~--- аутентификация на~основе риска (Risk-Based Authentication,
RBA). ACS больше не~делает жёсткий выбор \enquote{пароль или ничего},
а~оценивает риск и~выбирает подходящий сценарий.

\subsection*{Незаметный сценарий (\textit{frictionless flow})}

\highlight{Если ACS считает транзакцию низкорисковой, аутентификация происходит
\emph{невидимо} для~держателя карты}. 3DS Server отправляет контекст операции,
ACS оценивает его, возвращает положительный результат, и~транзакция передаётся
на~авторизацию без~дополнительной паузы на~экране оплаты. Этот сценарий сделал 3DS
пригодным для массовой электронной коммерции~--- дополнительная проверка
осталась в~процессе, но перестала разрывать оформление заказа на~каждом шаге.

Качество решения напрямую зависит от~контекста, который передаёт 3DS
Server. ACS получает сведения об~устройстве, браузере или приложении, адресе
доставки, истории покупок у~конкретного продавца, характере самой операции
и~поведенческих признаках. По~этому набору он решает, похожа~ли текущая покупка
на~привычное действие держателя карты или требует дополнительного
подтверждения. Чем полнее и~точнее эти данные, тем чаще эмитент пропускает
операцию незаметно для клиента, и~тем больше выручки сохраняет продавец.

Полный цикл сообщений в~EMV 3DS устроен так. Ещё до~AReq браузер держателя
карты может быть незаметно перенаправлен на~\textit{3DS Method URL} ACS~---
служебный URL, вызываемый до~AReq в~скрытом iframe, в~котором ACS собирает
сигнатуру устройства (\textit{device fingerprinting}; механика и~сигналы~---
в~разделе~\ref{sec:fraud-fingerprint}) и~привязывает её к~идентификатору
будущей транзакции. Затем 3DS Server отправляет в~Directory Server
Authentication Request (AReq), DS маршрутизирует его в~нужный ACS и~возвращает
Authentication Response (ARes). В~незаметном сценарии на~ARes цикл
и~заканчивается. Если ACS запрашивает подтверждение, используется пара
Challenge Request~/ Challenge Response (CReq/CRes) между ACS и~браузером или
3DS SDK (клиентская библиотека для~нативного мобильного приложения, заменяющая
браузер в~потоке аутентификации), а~по~её завершении ACS отправляет в~DS итог
в~Results Request (RReq); DS ретранслирует его 3DS Server, который отвечает
Results Response (RRes). Итоговый статус аутентификации доставляет в~3DS Server
именно DS, а~не~ACS напрямую. Статус приходит в~поле transStatus:
\texttt{Y}~--- аутентификация пройдена, \texttt{N}~--- не~пройдена,
\texttt{A}~--- попытка аутентификации (\textit{attempted authentication}; см.~раздел~\ref{sec:3ds-liability}).

Конкретный набор полей AReq определяет качество этих данных. Среди десятков
полей спецификации EMVCo выделяются группы, которые напрямую влияют на~долю
\textit{frictionless}~\cite{emvco-3ds-whitepaper}:

\begin{itemize}
  \item \textbf{deviceChannel}, параметры браузера
    (\texttt{browserScreen\allowbreak Height}, \texttt{browser\allowbreak Language} и~др.)~---
    сигнатура устройства, ACS сопоставляет текущую сессию с~ранее
    виденными;
  \item \textbf{threeDSRequestorAuthenticationInd}~--- тип операции (оплата,
    рекуррент, добавление карты), задаёт порог риска;
  \item \textbf{acctInfo}~--- возраст аккаунта у~продавца, дата последнего
    изменения адреса и~количество покупок за~6~месяцев, то~есть прямые
    индикаторы доверия;
  \item \textbf{merchantRiskIndicator}~--- сроки доставки, тип товара
    (физический, цифровой, подписка), предзаказ;
  \item \textbf{shipAddrMatch} и~адресные поля~--- расхождение адреса доставки
    с~адресом биллинга, сигнал для запроса подтверждения;
  \item \textbf{threeDSRequestorPriorAuthenticationInfo}~--- ссылка
    на~предыдущую успешную аутентификацию у~того же продавца, снижающая
    порог для повторных покупок.
\end{itemize}

Чем больше релевантных полей с~данными о~риске передаёт продавец, тем выше у~него
доля аутентификаций в~незаметном сценарии~--- транзакций, которые ACS пропускает
без~вызова дополнительного подтверждения. Точная доля \textit{frictionless}
и~доля завершения (\textit{completion rate}) для сценариев с~подтверждением (\textit{challenge})
зависят от~сегмента ТСП, региона и~качества данных по~риску; публичные
агрегированные числа по~ним редки, а~маркетинговые материалы вендоров не~дают
общего правила, поэтому в~книге они не~приводятся.\footnote{Внутренние
  A/B-измерения у~крупных PSP стабильно показывают, что улучшение качества
  переданных в~3DS-сообщении полей даёт измеримый рост доли \textit{frictionless}
  и~совокупной конверсии страницы оплаты. Конкретные цифры зависят от~внутреннего
сравнительного измерения продавца; общего правила нет.} Каждый
дополнительный запрос подтверждения отсеивает часть покупателей, поэтому качество данных
в~3DS-сообщении напрямую влияет на~конверсию.

\subsection*{Сценарий дополнительного подтверждения (\textit{challenge flow})}

Если ACS обнаруживает повышенный риск~--- например, новое устройство или нетипичную
сумму,~--- он запускает запрос дополнительного подтверждения. В~отличие
от~первой версии, протокол поддерживает несколько методов подтверждения
в~цифровом канале. ACS инициирует Challenge Request (CReq) и~передаёт его
напрямую в~браузер или 3DS SDK; после ответа держателя карты Challenge Response
(CRes) возвращается к~ACS~\cite{emvco-3ds-whitepaper}.

Метод подтверждения зависит от~канала: одноразовый пароль по~SMS или
push-подтверждение с~биометрией в~мобильном банке. 3DS v2 использует
интерфейсы, которыми держатель карты пользуется каждый день, и~не~зависит
от~редкого пароля~\cite{emvco-3ds-whitepaper}.

\begin{figure}[htbp]
  \centering
  \includefiguresvg[width=\linewidth]{assets/figures/ch11-3ds/3ds-v2-flow}
  \caption{Поток 3DS v2: незаметный сценарий (AReq/ARes) и~сценарий с~подтверждением (CReq/CRes, RReq/RRes).}\label{fig:3ds-v2-flow}
\end{figure}

\FloatBarrier
\section{Развязанная аутентификация}\label{sec:3ds-decoupled}

3DS v2 ввёл ещё один сценарий~--- развязанную аутентификацию (\textit{decoupled
authentication}), то~есть подтверждение, которое происходит \emph{не~в~момент}
покупки и~\emph{не~в~контексте} страницы оплаты. Обычный сценарий
дополнительного подтверждения предполагает, что держатель карты находится рядом
с~устройством и~может немедленно ответить на~запрос. Устройство интернета
вещей без~экрана~--- например, умная колонка или подключённый автомобиль~---
инициирует покупку без~участия человека. В~рекуррентном или ином 3RI-сценарии
(3RI~--- 3DS Requestor Initiated, аутентификация, инициированная не~держателем
карты, а~продавцом или его 3DS Requestor) запрос тоже формируется без~присутствия
держателя карты. Для таких случаев предназначена развязанная
аутентификация~\cite{emvco-3ds-decoupled}.

В~отличие от~обычного сценария с~подтверждением ACS подтверждает использование
развязанной аутентификации в~сообщениях 3DS и~затем передаёт итог 3DS Server
в~Results Request (RReq). Для продавца это промежуточный статус: заказ нельзя считать
окончательно подтверждённым, пока не~получен результат
аутентификации~\cite{emvco-3ds-decoupled}.

После этого 3DS Server либо периодически запрашивает итоговый статус, либо получает
его по~предусмотренному в~интеграции асинхронному каналу. Параллельно эмитент
связывается с~держателем карты доступным способом: push-уведомлением или
SMS~\cite{emvco-3ds-decoupled}.

Спецификация задаёт максимальное время ожидания (\textit{decoupled max
timeout}), после которого аутентификация считается неудавшейся. Если она
не~завершена вовремя, продавец либо отменяет заказ, либо отправляет транзакцию
без~результата аутентификации, принимая повышенный
риск~\cite{emvco-3ds-decoupled}.

Сценарий развязанной аутентификации для держателя карты требует отдельного
проектирования. Типичная формулировка~--- \enquote{Мы отправили вам уведомление
в~приложении: подтвердите, пожалуйста, покупку, и~мы отгрузим заказ}~---
фиксирует заказ условно, а~не немедленно. Система управления заказами должна
обработать промежуточный статус, таймаут и~уведомление при~подтверждении или
отказе. Держателю карты нужно объяснить, что подтверждение завершается
в~банковском приложении~\cite{emvco-3ds-decoupled}.

\section{Перенос ответственности}\label{sec:3ds-liability}

Перенос ответственности (\textit{liability shift}) означает договорное
перераспределение риска chargeback по~причине мошенничества с~продавца и~эквайера
на~эмитента. Он остаётся ключевым экономическим стимулом для внедрения 3DS, но
\highlight{включается не~автоматически}: зависит от~результата аутентификации,
статуса операции, правил сети и, в~Европе, от~того, использовалось~ли
освобождение от~SCA~\cite{visa-core-rules}. Различают четыре сценария:

\begin{itemize}
  \item \textbf{3DS выполнен, эмитент одобрил}~--- незаметный сценарий или
    подтверждение пройдено. При выполнении условий программы сети перенос
    ответственности работает, риск переходит к~эмитенту. Продавец защищён
    от~chargeback по~коду причины, связанному с~мошенничеством.
  \item \textbf{получен результат попытки аутентификации}~--- случай, когда
    ACS эмитента не~отвечает или эмитент
    не~участвует в~программе; результат фиксирует сервер сети в~режиме
    \textit{stand-in}, отвечающий от~имени эмитента. Исторически некоторые
    программы давали продавцу защиту и~в~этом случае, но универсальным
    правилом EMV 3DS это не~является: эффект зависит от~конкретной сети и~её
    текущей программы.
  \item \textbf{продавец не~использует 3DS}~--- перенос ответственности
    не~работает, продавец и~эквайер несут ответственность за~мошенничество.
  \item \textbf{продавец использует освобождение}~--- см.~следующий раздел.
    Перенос ответственности здесь \emph{не~работает}~--- продавец и~эквайер
    осознанно отказались от~аутентификации и~приняли риск на~себя.
\end{itemize}

Правила различаются по~сетям и~регионам: у~Visa и~Mastercard свои условия для
попытки аутентификации, рекуррентных операций и~отдельных категорий продавцов.
Наиболее изменчивы правила попытки аутентификации.

По~Visa Secure Program Guide эмитент в~программе на~EMV 3DS
возвращает \enquote{Not Authenticated}; ответ \enquote{Attempts}
за~неучаствующих эмитентов формирует Visa Attempts Server в~режиме
\textit{stand-in}~\cite{visa-secure-rules-update}.

В~России, где PSD2 не~действует, перенос ответственности для внутрироссийских
операций определяется правилами НСПК для карт Мир. До~марта 2022~года
аналогичные правила действовали для Visa и~Mastercard; после приостановки
операций этих сетей в~России к~внутрироссийским транзакциям они уже
не~применимы.

\section{Освобождения PSD2 и~специфика России}\label{sec:3ds-exemptions}

PSD2 и~регуляторные технические стандарты (Regulatory Technical Standards, RTS)
Европейской банковской службы (European Banking Authority,
EBA)~\cite{eba-rts-sca} требуют SCA для электронных платежей, кроме ряда
освобождений: \textit{low value} (операции на~небольшие суммы), TRA (Transaction Risk
Analysis~--- анализ риска транзакции эквайером или эмитентом), \textit{trusted
beneficiary} (доверенный получатель в~белом списке держателя карты), \textit{recurring}
(рекуррентные списания с~одинаковой суммой) и~\textit{secure corporate} (защищённые
корпоративные платежи по~выделенным протоколам). Общая модель и~состав
исключений разобраны в~разделе~\ref{sec:psd2-sca}.

Освобождение запрашивает продавец или эквайер, а~решение принимает эмитент.
Применение TRA или \textit{low value} не~отменяет его модели риска, и~если
транзакция выглядит подозрительно, эмитент вправе потребовать \textit{challenge}.
Использование освобождения может ослабить или полностью выключить перенос
ответственности: убирая дополнительную проверку, продавец принимает риск мошенничества
на~себя~\cite{eba-rts-sca}. 3DS v2 нужен и~в~этом случае как канал для
индикаторов риска, запрошенных освобождений и~итогового решения
эмитента~\cite{eba-rts-sca}.

\subsection*{Специфика России}

Россия не~входит в~зону действия PSD2, и~национальное законодательство
не~содержит прямого аналога PSD2/RTS с~тем же обязательным набором освобождений
и~их единым описанием. Необходимость аутентификации по~3DS для внутрироссийских
транзакций остаётся предметом политики эмитента и~правил карточной сети.

\section{3DS и~токенизация}\label{sec:3ds-tokenization}

Токенизация (глава~\ref{ch:tokenization}) и~3DS решают разные задачи. Токен
отвечает на~вопрос, какие именно карточные реквизиты участвуют в~операции;
3DS~--- действительно~ли её инициирует законный держатель карты. Вместе они
дают эмитенту более полную картину риска. В~сценарии \textit{card-on-file}
в~электронной коммерции с~EMV 3DS продавец передаёт в~3DS Server сетевой токен и~связанные
с~ним атрибуты; Directory Server запрашивает у~TSP (Token Service Provider,
поставщик услуг токенизации) детокенизацию с~верификацией (\textit{de-tokenisation with verification})~--- проверкой криптограммы и~доменных
ограничений на~стороне TSP~--- и~передаёт нужные данные в~ACS, а~ACS использует
их, чтобы решить, нужен~ли \textit{challenge} и~можно~ли сформировать значение аутентификации (\textit{authentication
  value}~--- криптографическое доказательство успешной аутентификации: CAVV
  у~Visa, AAV у~Mastercard, передаваемое
в~авторизации)~\cite{emvco-payment-tokenisation-guide}.

Для рекуррентных операций механика выглядит иначе. В~EMV 3DS версии 2.2 и~выше
последующие аутентификации повторных списаний \textit{recurring} или рассрочки
\textit{instalment} могут выполняться как 3RI-запросы. 3DS Requestor указывает
3RI Indicator и~передаёт DS Transaction~ID и/или ACS Transaction~ID (уникальные
  идентификаторы шага DS и~ACS, выработанные при~исходной клиентской
аутентификации), чтобы ACS мог соотнести новую операцию с~исходной
аутентификацией держателя карты~\cite{emvco-3ds-recurring}. В~руководстве EMVCo
по~токенизации указано, что в~сценариях MIT (Merchant-Initiated Transaction~---
  операция, инициированная продавцом по~ранее установленному мандату держателя
карты) свойства токена, способ его проверки
и~доменные ограничения наследуют контекст исходной клиентской
операции~\cite{emvco-payment-tokenisation-guide}. Подробно о~влиянии
токенизации и~3DS на~антифрод~--- в~главе~\ref{ch:fraud}.
